\section{Predicting a round}
\label{sec:do-the-checks-predict}

The three evaluation models yield \gpNovGainCases{} cases in total, of which
the outcome rule classifies \gpCases{}. We ask two questions of them. First,
does the frozen threshold rule predict what a single round will do before
training (Section~\ref{sec:main-result}), and what explains its errors
(Section~\ref{sec:errors})? Second, do Propositions~\ref{prop:novelty}
and~\ref{prop:agreement} hold quantitatively underneath the thresholds: does
novelty track the size of the gain (Section~\ref{sec:novelty-bound}), and does
split agreement isolate right answers without a task filter
(Section~\ref{sec:confirmed-right})?

\subsection{The main result}
\label{sec:main-result}
\label{sec:which-check}

Across the three held-out starting models, the prediction rule was correct in
\gpCorrect{} of \gpCases{} cases (\gpCorrectPct\%), compared with
\gpDidNothingPct\% for always predicting \emph{does nothing}, the most frequent
outcome (Table~\ref{tab:p9b-confusion}). Every prediction of \emph{pays off}
and of \emph{does nothing} was correct, all \gpHarmful{} harmful cases were
flagged as \emph{does harm}, and all nine errors were false alarms of harm,
meeting the pass mark set before evaluation. Broken down by starting model, the
rule was correct in \gpCorrectModelTwo{}, \gpCorrectModelThree{}, and
\gpCorrectModelFour{} of the \gpCasesModelTwo{} cases per seed.

% Generated by `uv run python -m gainpred.writeup`; do not edit by hand.
\begin{table}[t]
\centering
\caption{Predictions before training against outcomes after training, for the
  258 evaluation cases. Every error is a prediction of harm that did not occur.}
\label{tab:p9b-confusion}
\small
\begin{tabular}{lrrr}
\toprule
 & \multicolumn{3}{c}{Prediction} \\
\cmidrule(l){2-4}
Outcome & Pays off & Does nothing & Does harm \\
\midrule
Paid off & 31 & 0 & 6 \\
Did nothing & 0 & 175 & 3 \\
Did harm & 0 & 0 & 43 \\
\bottomrule
\end{tabular}
\end{table}


The \gpCases{} cases are not independent tests. They fall into eight kinds, and
within each kind the outcome is nearly always the same: what a round does is
largely determined by which teacher is used and whether the length is new
(Supplement, Table~\ref{S-tab:p9b-cases}). The checks are given neither piece
of information; they infer what the round will do from the model's answers
alone. The novelty check (including its 5\% retention floor) decided
\gpByNovelty{} cases, all correctly: lengths already learned, every round
taught by sample and vote, and rounds with no selection at lengths it had
already taught. Agreement decided \gpByAgreement{} cases, all correctly,
separating the teachers that paid off at new lengths from no selection, which
taught wrong answers there. The basics check decided the remaining
\gpByBasics{} cases and accounts for all nine errors. Each check falls into two
well-separated clusters with the threshold inside the gap between them
(Supplement, Figure~\ref{S-fig:checks}). Running the checks took
\gpChecksCostAdditionPct\% of a round's compute; starting the round and extrapolating
from its first 200 updates was less accurate, took \gpBaselineCostAdditionPct\% of a
round's compute, and required correct answers
(Supplement~\ref{S-app:prediction}).

\subsection{Novelty bounds the change}
\label{sec:novelty-bound}

The prediction rule uses novelty only as a cutoff, whereas
Proposition~\ref{prop:novelty} makes a quantitative statement: novelty bounds
the magnitude of the change a round can produce, and equals the gain when the
teacher is right. We tested this after all experiments on the same cases, with
no prediction recorded in advance. In \gpNovGainWithin{} of \gpNovGainCases{}
cases, test accuracy changed by no more than the novelty plus 3 points, the
margin allowed for a round that did nothing (Appendix~\ref{app:evidence},
Figure~\ref{fig:novelty-gain}). When the teacher is right, the proposition
predicts that novelty, pre-round student error, and accuracy gain coincide. On
fully trained starting models, the no-carry rule and the agreement rule are
right almost without exception: across their \gpNovGainSound{} cases, the gain
lay within 5 points of the novelty in \gpNovGainSoundChange{} (largest gap
\gpNovGainSoundChangeGap), including the \gpNovGainSoundBetween{} partially
learned cases where novelty sits away from 0\% and 100\% and all three numbers
move together. For a teacher known to be right, pre-round novelty is therefore
a direct forecast of the points a round will add. Without a check on teacher
accuracy, novelty remains only an upper bound: where sample and vote kept
answers, they mostly differed from the student's, yet accuracy never changed by
more than \gpNovGainVoteMax{} points.

\subsection{Confirmed answers are right}
\label{sec:confirmed-right}
\label{sec:agreement-filter}

Proposition~\ref{prop:agreement} requires $\varepsilon$ to be small for the
agreement check ($a_T \ge c - \varepsilon$) and small relative to $c$ for the
agreement rule ($1 - \varepsilon/c$), and $\varepsilon$ is the one quantity a
model cannot measure for itself. After all experiments finished, we measured it
with correct answers; before any answer was scored, we recorded one prediction
for each condition across the \gpTaFullCases{} split-based cases from fully
trained starting models: that $\varepsilon$ would be at most 2\% in every case,
and that at least 98\% of confirmed answers ($1 - \varepsilon/c$) would be
correct in every case. Both held. The largest $\varepsilon$ was
\gpTaPOneWorst\%, and the lowest accuracy among confirmed answers was
\gpTaPTwoWorst\% (Appendix~\ref{app:evidence},
Table~\ref{tab:teacher-accuracy}).

Agreement therefore identifies the right answers even of a teacher that is
mostly wrong. With no selection, the teacher was right on only
\gpTaNoSelectAccLow--\gpTaNoSelectAccHigh\% of its answers, yet at least
\gpTaNoSelectAccConfirmedLow\% of those confirmed by another split were
correct. This is why the agreement rule can replace a filter written for the
task: it encodes no knowledge of carrying, and after four rounds its student
was within \gpGapMean{} points of the no-carry student on average across 8 to
15 digits, at the price of taking \gpBuildRatio{} times as long to fill the
same training quota (Supplement~\ref{S-app:prediction}). The guarantee has a
limit. On weak starting models, which answer the pieces themselves poorly,
$\varepsilon$ reached \gpTaWeakWrongConfirmedHigh\% while $c$ also fell,
pulling confirmed accuracy down with it (Table~\ref{tab:teacher-accuracy}):
split agreement tests the composition step directly, and piece accuracy only
indirectly.


\subsection{Where the prediction was wrong}
\label{sec:errors}

All nine errors on classified cases come from the ten cases whose starting
model failed the basics check: the rule predicted harm in all ten and was right
once, because a model just below the 98\% pass mark still yields a teacher that
is mostly right. Without the basics check, the other two checks alone would
have been correct in \gpCorrectNoBasics{} of \gpCases{} cases
(\gpCorrectNoBasicsPct\%); we computed this afterwards, and the frozen rule
remains our main result.

The outcome rule leaves \gpNoOutcome{} further cases unclassified. In ten of
them, a length the model already answered correctly more than half the time
rose to 98\% or higher, while the rule predicted \emph{does nothing}. These are
errors in substance: all ten were taught by an accurate teacher, so novelty
equaled the student's remaining error, which lay below 50\% and was read by the
rule as nothing left to learn. The fault lies in the coarse 50\% cutoff rather
than in novelty itself, which came within \gpNovGainSoundChangeGap{} points of
the actual gain in every one of the ten cases. Counting all \gpNoOutcome{}
unclassified cases as errors gives \gpCorrectAllPct\%.

\paragraph{From single rounds to a loop.}
Taken together, these results separate the two quantities from the thresholds
wrapped around them. The frozen rule predicts single-round outcomes in
\gpCorrectPct\% of classified cases and flags every harmful round, and its
misses---false alarms from the 98\% basics floor and partially learned lengths
missed by the 50\% novelty cutoff---come from discretizing continuous signals
rather than from failures of the bounds themselves. Underneath those cutoffs,
split agreement isolates right answers ($\ge\gpTaPTwoWorst\%$ accurate on fully
trained models) even when an unfiltered teacher is mostly wrong, and novelty
forecasts the resulting gain within a few points whenever the teacher is right.
Those are the two pieces a multi-round loop needs: a task-free filter to keep
composition errors out of training, and a quantitative measure of whether
anything remains to be learned.
